\section{Билет 29. Поток и плотность потока энергии. Вектор Умова. Интенсивность волны.}

\begin{definition}
\textbf{Потоком энергии} волны через некоторую поверхность называется энергия, переносимая волной через эту поверхность за единицу времени. Единица измерения потока энергии в СИ --- ватт (Вт).
\end{definition}

\begin{definition}
\textbf{Вектор плотности потока энергии} (вектор Умова) --- векторная величина, характеризующая направление и плотность переноса энергии волной. Введён русским физиком Н.А. Умовым (1846--1915).
\begin{itemize}
    \item \textbf{Направление} вектора Умова $\vec{S}$ совпадает с направлением распространения волны (нормалью к волновой поверхности).
    \item \textbf{Модуль} вектора равен энергии, переносимой волной за единицу времени через единичную площадку, расположенную перпендикулярно направлению распространения:
    \begin{equation}
        S = \frac{dE}{dt \, dS_{\perp}}. \label{eq:umov_def}
    \end{equation}
\end{itemize}
Единица измерения --- Вт/м².
\end{definition}

\begin{theorem}[Выражение для вектора Умова упругой волны]
Рассмотрим плоскую монохроматическую волну:
\begin{equation}
    \xi(x,t) = A \cos(\omega t - kx).
\end{equation}
Полная плотность энергии волны (сумма кинетической и потенциальной) равна:
\begin{equation}
    w(x,t) = \rho A^2 \omega^2 \sin^2(\omega t - kx). \label{eq:energy_density_29}
\end{equation}
Энергия переносится волной со скоростью $v$. За время $dt$ через площадку $dS_{\perp}$ проходит энергия, содержащаяся в объёме $dV = v \, dt \, dS_{\perp}$:
\begin{equation}
    dE = w \, dV = w \, v \, dt \, dS_{\perp}.
\end{equation}
Подставляя в определение \eqref{eq:umov_def}, получаем мгновенное значение модуля вектора Умова:
\begin{equation}
    S(t) = w(t) \cdot v = \rho v A^2 \omega^2 \sin^2(\omega t - kx). \label{eq:umov_instant}
\end{equation}
\end{theorem}

\begin{definition}
\textbf{Интенсивностью волны} $I$ называется среднее по времени значение модуля вектора плотности потока энергии:
\begin{equation}
    I = \langle S \rangle = \langle w \rangle \cdot v.
\end{equation}
Учитывая, что среднее значение $\sin^2$ за период равно $1/2$, из \eqref{eq:umov_instant} получаем формулу для интенсивности плоской монохроматической волны:
\begin{equation}
    I = \frac{1}{2} \rho v A^2 \omega^2. \label{eq:intensity_plane}
\end{equation}
Интенсивность пропорциональна квадрату амплитуды и квадрату частоты волны.
\end{definition}

\begin{remark}
\textbf{Сферическая волна.} Если волна распространяется от точечного источника в непоглощающей среде, то полная мощность $P$, излучаемая источником, распределяется по сферической поверхности радиуса $r$:
\begin{equation}
    P = I(r) \cdot 4\pi r^2 = \text{const}.
\end{equation}
Отсюда интенсивность и амплитуда сферической волны убывают с расстоянием по законам:
\begin{equation}
    I(r) = \frac{I_0 r_0^2}{r^2}, \quad A(r) = \frac{A_0 r_0}{r}. \label{eq:spherical_decay}
\end{equation}
Это объясняет, почему звук от точечного источника становится тише по мере удаления.
\end{remark}
